\documentclass[11pt,a4paper]{article}

\usepackage[margin=2.4cm]{geometry}
\usepackage{amsmath,amssymb,amsfonts}
\usepackage{braket}
\usepackage{booktabs}
\usepackage{longtable}
\usepackage{array}
\usepackage{xcolor}
\usepackage{listings}
\usepackage{microtype}
\usepackage[colorlinks=true,linkcolor=blue!50!black,citecolor=blue!50!black,urlcolor=blue!50!black]{hyperref}

\lstset{basicstyle=\small\ttfamily,
        keywordstyle=\color{blue!60!black},
        commentstyle=\color{green!40!black},
        stringstyle=\color{purple!60!black},
        language=Matlab,
        breaklines=true,
        columns=fullflexible,
        keepspaces=true,
        frame=single,
        rulecolor=\color{black!30},
        xleftmargin=1em,
        xrightmargin=1em}

\newcommand{\code}[1]{\texttt{#1}}
\newcommand{\Tr}{\operatorname{Tr}}
\newcommand{\op}[1]{\hat{#1}}
\newcommand{\sop}[1]{\hat{\hat{#1}}}

\title{\textbf{The Spinach Optimal Control Module}\\[2mm]
       \Large User Manual}
\author{Talos\thanks{AI research assistant, Weizmann Institute of Science.
        The module itself is the work of David Goodwin, Uluk Rasulov,
        Marcel Keitel, Ilya Kuprov, and the Spinach developer team.}}
\date{27 August 2026 \\[1mm] \normalsize Spinach 2.12, optimal control
      module architecture of pull request \#332}

\begin{document}

\maketitle

\tableofcontents
\clearpage

%===============================================================
\section{Physics and mathematics of GRAPE}
%===============================================================

\subsection{The quantum optimal control problem}

Magnetic resonance systems evolve under an equation of motion that Spinach
always writes in the form
\begin{equation}
  \frac{d}{dt}\,\boldsymbol{\rho}(t)
  =-i\,\mathbf{L}(t)\,\boldsymbol{\rho}(t),
  \qquad
  \boldsymbol{\rho}(t)
  =\exp_{(\mathrm{o})}\!\left[-i\!\int_0^t \mathbf{L}(\tau)\,d\tau\right]
   \boldsymbol{\rho}(0),
  \label{eq:eom}
\end{equation}
where $\exp_{(\mathrm{o})}$ denotes a time-ordered exponential. In Hilbert
space $\boldsymbol{\rho}$ is a density matrix and $\mathbf{L}$ is the
commutation superoperator with a Hermitian Hamiltonian, so that the
propagation over an interval of constant $\mathbf{H}$ is a unitary sandwich
\begin{equation}
  \boldsymbol{\rho}(t+\Delta t)
  =\mathbf{P}\,\boldsymbol{\rho}(t)\,\mathbf{P}^{\dagger},
  \qquad
  \mathbf{P}=\exp(-i\mathbf{H}\Delta t).
  \label{eq:sandwich}
\end{equation}
In Liouville space $\boldsymbol{\rho}$ is a state vector and $\mathbf{L}$
is a Liouvillian that may contain relaxation and chemical kinetics; it is
then not Hermitian, the propagation
\begin{equation}
  \boldsymbol{\rho}(t+\Delta t)=\exp(-i\mathbf{L}\Delta t)\,
  \boldsymbol{\rho}(t)
  \label{eq:liouvprop}
\end{equation}
is not unitary, and dissipative dynamics is handled without further
ceremony. This is the principal reason why the workhorse formalism of the
optimal control module is Liouville space.

The control problem is stated as follows. The evolution generator is split
into an uncontrollable \emph{drift} and a linear combination of
\emph{control operators} $\{\mathbf{C}_k\}$, $k=1,\ldots,K$, whose
coefficients are at the disposal of the instrument:
\begin{equation}
  \mathbf{L}(t)=\mathbf{L}_0(t)+\sum_{k=1}^{K}c_k(t)\,\mathbf{C}_k.
  \label{eq:splitting}
\end{equation}
Typical control operators are $\op{L}_X$ and $\op{L}_Y$ magnetisation
operators on the channels of the instrument; the drift contains chemical
shifts, couplings, exchange, and relaxation. The pulse of total duration
$T$ is discretised over $N$ intervals $\Delta t_1,\ldots,\Delta t_N$. Two
discretisations of $c_k(t)$ are supported:
\begin{description}
  \item[Piecewise-constant (\code{rectangle}).] The coefficient
  $c_k(t)=c_{kn}$ is constant on the $n$-th interval; the waveform array
  has $N$ columns. This is the classical GRAPE setting of Khaneja \emph{et
  al.}~\cite{khaneja2005}.
  \item[Piecewise-linear (\code{trapezium}).] The coefficient is a linear
  interpolation between \emph{nodal} values $c_{kn}$ specified at the
  $N{+}1$ interval edges; the waveform array has $N{+}1$ columns. This
  matches what waveform generators actually emit, and is integrated with a
  product quadrature (Section~\ref{sec:trapezium})~\cite{rasulov2023}.
\end{description}
In Spinach the user-facing waveform is dimensionless and order unity: the
physical coefficient is $c_{kn}=p\,u_{kn}$, where $p$ is the nominal
control power (rad/s) supplied in \code{control.pwr\_levels}, and
$u_{kn}\in[-1,1]$ is the normalised waveform that the optimiser
manipulates. All derivatives reported to the optimiser are with respect to
$u_{kn}$; conversions are handled internally by the chain rule.

Starting from an initial state $\boldsymbol{\rho}_0$, the pulse should
arrive as closely as possible at a target $\boldsymbol{\sigma}$. With the
overlap
\begin{equation}
  o=\braket{\boldsymbol{\sigma}|\boldsymbol{\rho}(T)}
  =\begin{cases}
     \boldsymbol{\sigma}^{\dagger}\boldsymbol{\rho}(T)
       & \text{(Liouville space)}\\[1mm]
     \Tr\!\left[\boldsymbol{\sigma}^{\dagger}\boldsymbol{\rho}(T)\right]
       & \text{(Hilbert space)}
   \end{cases}
  \label{eq:overlap}
\end{equation}
three fidelity functionals are supported
(\code{control.fidelity}):
\begin{equation}
  f_{\mathrm{real}}=\operatorname{Re}o,\qquad
  f_{\mathrm{imag}}=\operatorname{Im}o,\qquad
  f_{\mathrm{square}}=|o|^{2}.
  \label{eq:fidelities}
\end{equation}
When states are normalised, $f_{\mathrm{real}}\in[-1,+1]$ and
$f_{\mathrm{square}}\in[0,1]$; the square functional is insensitive to the
overall phase of the final state. The optimiser maximises the total
objective
\begin{equation}
  F(\mathbf{u})=f(\mathbf{u})-\sum_{j} w_j\,P_j(\mathbf{u}),
  \label{eq:objective}
\end{equation}
where $P_j$ are penalty functionals (Section~\ref{sec:penalties}) with
non-negative weights $w_j$ that hold the waveform inside instrumental
constraints. In ensemble optimisations (Section~\ref{sec:ensembles}) the
fidelity is additionally averaged over a set of system instances.

GRAPE --- gradient ascent pulse engineering~\cite{khaneja2005} --- is the
observation that, for piecewise-defined waveforms, the gradient
$\partial f/\partial u_{kn}$ (and, with more effort, the Hessian) can be
computed \emph{analytically} at a cost comparable to that of the
simulation itself, after which any standard numerical maximiser converges
rapidly. Spinach implements first-order (LBFGS) and second-order (Newton,
with two Hessian assembly strategies) optimisers over both integrators and
both spaces; the derivations follow.

\subsection{Directional derivatives of the matrix exponential}
\label{sec:dirdiff}

Everything in GRAPE reduces to derivatives of interval propagators
$\mathbf{P}=\exp[-i(\mathbf{A}+\varepsilon\mathbf{B})\Delta t]$ with
respect to $\varepsilon$ at $\varepsilon=0$ --- \emph{directional}
derivatives of the matrix exponential in the direction $\mathbf{B}$.
Differentiating the exponential series term by term,
\begin{equation}
  \mathcal{D}_{\mathbf{B}}\exp(-i\mathbf{A}\Delta t)
  =\left.\frac{\partial}{\partial\varepsilon}
   e^{-i(\mathbf{A}+\varepsilon\mathbf{B})\Delta t}\right|_{\varepsilon=0}
  =-i\!\int_0^{\Delta t}
   e^{-i\mathbf{A}(\Delta t-\tau)}\,\mathbf{B}\,
   e^{-i\mathbf{A}\tau}\,d\tau,
  \label{eq:duhamel}
\end{equation}
a nested integral that would be expensive head-on. The auxiliary matrix
technique of Van Loan~\cite{vanloan1978} and Najfeld and
Havel~\cite{najfeld1995}, in the form of Eq.~16 of Goodwin and
Kuprov~\cite{goodwin2015}, obtains it from a single exponential of a
block-bidiagonal matrix:
\begin{equation}
  \exp\!\left(-i
  \begin{bmatrix}
    \mathbf{A} & \mathbf{B}\\
    \mathbf{0} & \mathbf{A}
  \end{bmatrix}\Delta t\right)
  =
  \begin{bmatrix}
    \mathbf{P} & \mathcal{D}_{\mathbf{B}}\mathbf{P}\\
    \mathbf{0} & \mathbf{P}
  \end{bmatrix},
  \qquad \mathbf{P}=e^{-i\mathbf{A}\Delta t}.
  \label{eq:auxmat2}
\end{equation}
The block structure telescopes: with $N$ diagonal blocks $\mathbf{A}$ and
directions $\mathbf{B}_1,\ldots,\mathbf{B}_{N-1}$ on the superdiagonal,
the top-right corner of the exponential contains the nested integral of
the ordered product of all the directions; for equal directions
$\mathbf{B}$ it equals the $n$-th derivative divided by $n!$. In
particular, the $3\times3$ block matrix
\begin{equation}
  \exp\!\left(-i
  \begin{bmatrix}
    \mathbf{A} & \mathbf{B}_1 & \mathbf{0}\\
    \mathbf{0} & \mathbf{A}   & \mathbf{B}_2\\
    \mathbf{0} & \mathbf{0}   & \mathbf{A}
  \end{bmatrix}\Delta t\right)
\end{equation}
delivers in its top-right block the mixed second directional derivative
(the ordered half; the full mixed derivative is the sum of both orderings,
and for $\mathbf{B}_1{=}\mathbf{B}_2{=}\mathbf{B}$ the block is
$\tfrac{1}{2}\mathcal{D}^2_{\mathbf{B}}\mathbf{P}$). The kernel utility
\code{dirdiff.m} implements exactly this: it builds the $N$-block
auxiliary matrix, exponentiates it with the standard Spinach
\code{propagator()} machinery at tightened tolerance, and returns
$\{\mathbf{P},\,\mathcal{D}\mathbf{P},\,\mathcal{D}^2\mathbf{P},\ldots\}$
with the factorial factors restored.

Two properties make this exact and cheap: the auxiliary matrix is highly
sparse when $\mathbf{A}$ and $\mathbf{B}$ are sparse, and --- crucially
for Liouville space --- the derivative is never needed as a matrix, only
its \emph{action on a state}. Acting with Eq.~\eqref{eq:auxmat2} on a
stacked vector,
\begin{equation}
  \exp\!\left(-i
  \begin{bmatrix}
    \mathbf{A} & \mathbf{B}\\
    \mathbf{0} & \mathbf{A}
  \end{bmatrix}\Delta t\right)
  \begin{bmatrix}\mathbf{0}\\ \boldsymbol{\rho}\end{bmatrix}
  =
  \begin{bmatrix}
    (\mathcal{D}_{\mathbf{B}}\mathbf{P})\,\boldsymbol{\rho}\\
    \mathbf{P}\,\boldsymbol{\rho}
  \end{bmatrix},
  \label{eq:auxvec}
\end{equation}
so one call to the Spinach matrix-exponential-times-vector routine
\code{step()} on a doubled-dimension sparse generator returns the
derivative applied to the current state. This \emph{auxiliary vector
trick} is the computational core of the Liouville-space GRAPE gradient.

\subsection{GRAPE in Liouville space}
\label{sec:grapeliouv}

Let the piecewise-constant generator on interval $n$ be
\begin{equation}
  \mathbf{L}_n=\mathbf{L}_0^{(n)}+\sum_{k=1}^{K}c_{kn}\mathbf{C}_k,
  \qquad
  \mathbf{P}_n=\exp(-i\mathbf{L}_n\Delta t_n),
\end{equation}
where the drift $\mathbf{L}_0^{(n)}$ cycles through the user-supplied
drift array (one drift, or one per time slice for time-dependent drifts,
e.g.\ under magic angle spinning). The final state is
$\boldsymbol{\rho}(T)=\mathbf{P}_N\cdots\mathbf{P}_1\boldsymbol{\rho}_0$
and the overlap is
$o=\boldsymbol{\sigma}^{\dagger}\mathbf{P}_N\cdots\mathbf{P}_1
\boldsymbol{\rho}_0$. Differentiating with respect to $c_{kn}$ touches
only $\mathbf{P}_n$:
\begin{equation}
  \frac{\partial o}{\partial c_{kn}}
  =\boldsymbol{\sigma}^{\dagger}\,
   \mathbf{P}_N\cdots\mathbf{P}_{n+1}
   \left(\mathcal{D}_{\mathbf{C}_k}\mathbf{P}_n\right)
   \mathbf{P}_{n-1}\cdots\mathbf{P}_1\boldsymbol{\rho}_0
  =\boldsymbol{\chi}_{n}^{\dagger}
   \left(\mathcal{D}_{\mathbf{C}_k}\mathbf{P}_n\right)
   \boldsymbol{\rho}_{n-1},
  \label{eq:gradliouv}
\end{equation}
where the \emph{forward trajectory} and the \emph{adjoint (backward)
trajectory} are
\begin{equation}
  \boldsymbol{\rho}_{n}=\mathbf{P}_n\cdots\mathbf{P}_1
  \boldsymbol{\rho}_0,
  \qquad
  \boldsymbol{\chi}_{n}
  =\mathbf{P}_{n+1}^{\dagger}\cdots\mathbf{P}_N^{\dagger}\,
  \boldsymbol{\sigma}.
  \label{eq:trajectories}
\end{equation}
Both trajectories cost one propagation sweep each; all $KN$ gradient
elements then follow from Eq.~\eqref{eq:auxvec} at one doubled-dimension
\code{step()} call each, with the stacked vector
$[\,\mathbf{0};\boldsymbol{\rho}_{n-1}]$ and the auxiliary generator
built from $\mathbf{L}_n$ and $\mathbf{C}_k$. No propagator and no
derivative is ever stored as a matrix.

Because Liouvillians need not be Hermitian,
$\mathbf{P}_n^{\dagger}\neq\mathbf{P}_n^{-1}$ in general, and the adjoint
trajectory is generated by propagating with the \emph{adjoint generator
over negative time}:
$\mathbf{P}_n^{\dagger}=\exp(-i\mathbf{L}_n^{\dagger}(-\Delta t_n))$. The
kernel therefore builds backward generators from adjoint drifts and
adjoint control operators and steps them with $-\Delta t$; for Hermitian
controls this is invisible, and for deliberately non-Hermitian
(dissipative) controls it produces the correct adjoint. The same
convention gives correct gradients in the presence of relaxation, where
fidelity is genuinely lost along the way.

\emph{Keyholes} --- user-supplied linear idempotent maps (projectors)
applied to the state between intervals --- commute with this scheme: the
forward pass applies the keyhole after each step, the backward pass
applies it before the adjoint step at the mirrored position, and since
projectors are Hermitian the gradient formula~\eqref{eq:gradliouv} is
unchanged. Frozen waveform elements (a logical mask in
\code{control.freeze}) are simply skipped, and their gradient entries set
to zero.

\subsection{GRAPE in Hilbert space}
\label{sec:grapehilb}

In the \code{zeeman-hilb} formalism the state is a density matrix, the
propagation is the unitary sandwich of Eq.~\eqref{eq:sandwich}, and all
generators are required to be Hermitian (this is enforced at problem
setup). The trajectories are
\begin{equation}
  \boldsymbol{\rho}_{n}
  =\mathbf{P}_n\boldsymbol{\rho}_{n-1}\mathbf{P}_n^{\dagger},
  \qquad
  \boldsymbol{\chi}_{n}
  =\mathbf{P}_{n+1}^{\dagger}\boldsymbol{\chi}_{n+1}\mathbf{P}_{n+1},
  \qquad \boldsymbol{\chi}_{N}=\boldsymbol{\sigma},
\end{equation}
the overlap is the Frobenius inner product
$o=\Tr[\boldsymbol{\sigma}^{\dagger}\boldsymbol{\rho}_N]$, and the product
rule applied to the sandwich gives
\begin{equation}
  \frac{\partial o}{\partial c_{kn}}
  =\Braket{\boldsymbol{\chi}_{n}\,|\,
   (\mathcal{D}_{\mathbf{C}_k}\mathbf{P}_n)\,
   \boldsymbol{\rho}_{n-1}\mathbf{P}_n^{\dagger}
   +\mathbf{P}_n\boldsymbol{\rho}_{n-1}\,
   (\mathcal{D}_{\mathbf{C}_k}\mathbf{P}_n)^{\dagger}},
  \label{eq:gradhilb}
\end{equation}
with $\braket{\mathbf{A}|\mathbf{B}}=\Tr[\mathbf{A}^{\dagger}\mathbf{B}]$.
Here matrix dimensions are those of Hilbert, not Liouville, space, so
propagators and their derivatives are computed and stored as matrices:
\code{dirdiff.m} with two blocks returns $\{\mathbf{P}_n,
\mathcal{D}\mathbf{P}_n\}$ in one auxiliary exponential per $(k,n)$ pair.
Second derivatives for the Newton optimiser use the three-block auxiliary
matrix for $\mathcal{D}^2\mathbf{P}$ (equal and mixed directions) and the
second-order product rule for the sandwich,
\begin{equation}
  \partial^2(\mathbf{P}\boldsymbol{\rho}\mathbf{P}^{\dagger})
  =(\mathcal{D}^2\mathbf{P})\boldsymbol{\rho}\mathbf{P}^{\dagger}
  +2\,(\mathcal{D}\mathbf{P})\boldsymbol{\rho}
      (\mathcal{D}\mathbf{P})^{\dagger}
  +\mathbf{P}\boldsymbol{\rho}(\mathcal{D}^2\mathbf{P})^{\dagger},
\end{equation}
distributed over interval pairs in the same way as in Liouville space
(Section~\ref{sec:hessians}).

\subsection{Fidelity functionals and their derivatives}
\label{sec:fidelityfun}

All propagator differentiation above concerns the complex overlap $o$ and
its derivatives $g_{kn}=\partial o/\partial c_{kn}$ and
$h_{kn,jm}=\partial^2 o/\partial c_{kn}\partial c_{jm}$. The requested
fidelity functional is applied at the very end:
\begin{align}
  f=\operatorname{Re}o:&\quad
    \nabla f=\operatorname{Re}\mathbf{g},\quad
    \nabla^2 f=\operatorname{Re}\mathbf{h};\\
  f=\operatorname{Im}o:&\quad
    \nabla f=\operatorname{Im}\mathbf{g},\quad
    \nabla^2 f=\operatorname{Im}\mathbf{h};\\
  f=|o|^2=o\bar{o}:&\quad
    \nabla f=\bar{o}\,\mathbf{g}+o\,\bar{\mathbf{g}}
    =2\operatorname{Re}(\bar{o}\,\mathbf{g}),\\
  &\quad
    \nabla^2 f=\bar{o}\,\mathbf{h}+o\,\bar{\mathbf{h}}
    +\mathbf{g}\bar{\mathbf{g}}^{\mathrm{T}}
    +\bar{\mathbf{g}}\mathbf{g}^{\mathrm{T}}.
\end{align}
The \code{real} functional is the default and is appropriate for
state-to-state transfer with a fixed target phase; \code{square} removes
the phase sensitivity; \code{imag} is occasionally useful for quadrature
targets.

\subsection{The piecewise-linear (trapezium) integrator}
\label{sec:trapezium}

Real instruments emit waveforms that ramp linearly between the points of
the digital-to-analogue converter. On interval $n$ the generator
interpolates linearly between its edge values,
$\mathbf{G}(t)=\mathbf{G}_L+(\mathbf{G}_R-\mathbf{G}_L)\,t/\Delta t$,
where $\mathbf{G}_{L}$ and $\mathbf{G}_{R}$ are built from the nodal
control coefficients $c_{k,n}$ and $c_{k,n+1}$. The exact propagator is
the time-ordered exponential; expanding it in the Magnus series and
keeping the terms through second order gives
\begin{align}
  \boldsymbol{\Omega}
  &=-i\!\int_0^{\Delta t}\!\mathbf{G}(t)\,dt
    +\frac{1}{2}\int_0^{\Delta t}\!\!\!\int_0^{t_1}
    \left[-i\mathbf{G}(t_1),-i\mathbf{G}(t_2)\right]dt_2\,dt_1
    +\cdots\notag\\
  &=-i\Delta t\,\frac{\mathbf{G}_L+\mathbf{G}_R}{2}
    +\frac{\Delta t^2}{12}\left[\mathbf{G}_L,\mathbf{G}_R\right]
    +O(\Delta t^{5}),
\end{align}
where the double integral was evaluated using
$[\mathbf{G}(t_1),\mathbf{G}(t_2)]
=\frac{t_2-t_1}{\Delta t}[\mathbf{G}_L,\mathbf{G}_R]$. The interval
propagator is therefore written as a standard exponential of an
\emph{effective generator},
\begin{equation}
  \mathbf{P}=\exp(-i\,\mathbf{G}_{\mathrm{eff}}\,\Delta t),
  \qquad
  \mathbf{G}_{\mathrm{eff}}
  =\frac{\mathbf{G}_L+\mathbf{G}_R}{2}
   +\frac{i\Delta t}{12}\left[\mathbf{G}_L,\mathbf{G}_R\right],
  \label{eq:inquad}
\end{equation}
the Iserles--N{\o}rsett product quadrature~\cite{iserles1999} used
throughout Spinach for piecewise-linear
propagation~\cite{rasulov2023}. The kernel routine \code{step()} applies
this quadrature when it is given the edge and midpoint generators of an
interval; the optimal control code supplies
$\{\mathbf{G}_L,(\mathbf{G}_L{+}\mathbf{G}_R)/2,\mathbf{G}_R\}$ for each
forward step and the adjoint triple in reverse order for each backward
step.

Derivatives require one extra chain-rule stage, because each nodal
coefficient now enters the effective generator both linearly and through
the commutator. Differentiating Eq.~\eqref{eq:inquad} with respect to the
left-edge and right-edge coefficients of control $k$:
\begin{align}
  \frac{\partial\mathbf{G}_{\mathrm{eff}}}{\partial c_{k}^{(L)}}
  &=\frac{\mathbf{C}_k}{2}
   +\frac{i\Delta t}{12}\left[\mathbf{C}_k,\mathbf{G}_R\right]
  =\frac{\mathbf{C}_k}{2}
   +\frac{i\Delta t}{12}\left(\left[\mathbf{C}_k,\mathbf{D}_R\right]
   +\sum_{m}c_{m}^{(R)}\left[\mathbf{C}_k,\mathbf{C}_m\right]\right),
   \label{eq:dgeffl}\\
  \frac{\partial\mathbf{G}_{\mathrm{eff}}}{\partial c_{k}^{(R)}}
  &=\frac{\mathbf{C}_k}{2}
   +\frac{i\Delta t}{12}\left[\mathbf{G}_L,\mathbf{C}_k\right]
  =\frac{\mathbf{C}_k}{2}
   +\frac{i\Delta t}{12}\left(\left[\mathbf{D}_L,\mathbf{C}_k\right]
   +\sum_{m}c_{m}^{(L)}\left[\mathbf{C}_m,\mathbf{C}_k\right]\right),
   \label{eq:dgeffr}
\end{align}
where $\mathbf{D}_{L,R}$ are the edge drifts. The control--control
commutators $[\mathbf{C}_k,\mathbf{C}_m]$ do not depend on the waveform;
they are computed once at problem setup, stored in
\code{control.cc\_comm}, and a logical table of exactly commuting pairs
(\code{control.cc\_comm\_idx}) skips the vanishing terms. The kernel
function \code{aux\_mat.m} assembles
Eqs.~\eqref{eq:dgeffl}--\eqref{eq:dgeffr} and packs them into the
auxiliary block matrices of Section~\ref{sec:dirdiff} with
$\mathbf{A}=\mathbf{G}_{\mathrm{eff}}$ and
$\mathbf{B}=\partial\mathbf{G}_{\mathrm{eff}}/\partial c$; the propagator
derivative then follows from the same auxiliary vector trick, applied per
interval edge.

Since every interior node is the right edge of one interval and the left
edge of the next, its gradient is a two-term product rule,
\begin{equation}
  \frac{\partial o}{\partial c_{kn}}
  =\boldsymbol{\chi}_{n}^{\dagger}
   \left(\frac{\partial\mathbf{P}_n}{\partial c_{kn}}\right)
   \boldsymbol{\rho}_{n-1}
  +\boldsymbol{\chi}_{n-1}^{\dagger}
   \left(\frac{\partial\mathbf{P}_{n-1}}{\partial c_{kn}}\right)
   \boldsymbol{\rho}_{n-2},
\end{equation}
with the first and the last node contributing one term only. The
trapezium integrator is available in both spaces; analytic Hessians are
not implemented for it, so the optimiser choices are the quasi-Newton
family (Section~\ref{sec:optimisers}).

\subsection{Analytic Hessians}
\label{sec:hessians}

For the piecewise-constant integrator, second derivatives of the overlap
come in two kinds. \emph{Same-interval} elements need the second
directional derivative of a single propagator,
\begin{equation}
  \frac{\partial^2 o}{\partial c_{kn}\partial c_{jn}}
  =\boldsymbol{\chi}_{n}^{\dagger}\,
   \mathcal{D}^{2}_{\mathbf{C}_k\mathbf{C}_j}\mathbf{P}_n\,
   \boldsymbol{\rho}_{n-1},
\end{equation}
obtained from the three-block auxiliary matrix of
Section~\ref{sec:dirdiff}; the code doubles the extracted ordered block,
which for $k=j$ is exactly $\mathcal{D}^2\mathbf{P}$ and for $k\neq j$
combines with the transposed assembly below into the symmetric mixed
derivative. \emph{Cross-interval} elements ($m<n$) factor into two first
derivatives separated by ordinary propagation,
\begin{equation}
  \frac{\partial^2 o}{\partial c_{kn}\partial c_{jm}}
  =\boldsymbol{\chi}_{n}^{\dagger}
   \left(\mathcal{D}_{\mathbf{C}_k}\mathbf{P}_n\right)
   \mathbf{P}_{n-1}\cdots\mathbf{P}_{m+1}
   \left(\mathcal{D}_{\mathbf{C}_j}\mathbf{P}_m\right)
   \boldsymbol{\rho}_{m-1}.
  \label{eq:hesscross}
\end{equation}
Two assembly strategies are implemented:
\begin{description}
\item[\code{newton}.] The derivative actions
$(\mathcal{D}_{\mathbf{C}_j}\mathbf{P}_m)\boldsymbol{\rho}_{m-1}$ are
stored for all $(j,m)$ during a single sweep; then, for each $n$, the
left factor
$\boldsymbol{\chi}_{n}^{\dagger}(\mathcal{D}_{\mathbf{C}_k}\mathbf{P}_n)$
is propagated \emph{backwards} interval by interval with the adjoint
generators, forming inner products with the stored right factors along
the way. The cost is one backward sweep per $(k,n)$ pair, all of it
matrix-vector.
\item[\code{goodwin}.] Goodwin's acceleration~\cite{goodwin2016} projects
every derivative action to the common time origin using cumulative
propagators
$\mathbf{Q}_n=\mathbf{P}_n\cdots\mathbf{P}_1$:
Eq.~\eqref{eq:hesscross} becomes a single inner product between
$\mathbf{Q}_{n-1}$-projected left rows and $\mathbf{Q}$-projected right
columns, at the price of computing and storing the cumulative propagators
as matrices. It trades memory for speed on problems where the propagators
are affordable; it also requires unitary propagation, so all generators
must be Hermitian (checked at setup).
\end{description}
The assembled Hessian is symmetrised, the rows and columns of frozen
waveform elements are cleared with unit diagonal, and the fidelity chain
rule of Section~\ref{sec:fidelityfun} is applied. Analytic Hessians are
available in both spaces for the rectangle integrator in ordinary
(non-steady-state, undistorted) optimisations.

\subsection{Ensemble optimal control}
\label{sec:ensembles}

A pulse is useful when it works over the instrument and sample
inhomogeneities. Spinach optimises the \emph{ensemble average fidelity}
\begin{equation}
  \bar{f}=\frac{1}{M}\sum_{i=1}^{M}f_i,
  \qquad
  \nabla\bar{f}=\frac{1}{M}\sum_{i=1}^{M}\nabla f_i,
\end{equation}
where each ensemble member $i$ is one fully specified simulation --- a
\emph{case}. Cases are enumerated as the Cartesian product of six
ensemble dimensions:
\begin{center}
\begin{tabular}{ll}
\toprule
dimension & source \\
\midrule
initial--target state pair      & \code{control.rho\_init} / \code{rho\_targ} cells\\
drift generator                 & \code{control.drifts} cell array\\
control power level             & \code{control.pwr\_levels} vector ($B_1$ miscalibration)\\
resonance offset combination    & \code{control.offsets} / \code{off\_ops} ($B_0$ distribution)\\
phase cycle line                & \code{control.phase\_cycle} rows\\
distortion function set         & \code{control.distortion} rows\\
\bottomrule
\end{tabular}
\end{center}
Offsets enter by adding $2\pi\,\delta_{\nu}\,\mathbf{O}_c$ to the drift of
the case, where $\delta_\nu$ is the offset value in Hz on channel $c$ and
$\mathbf{O}_c$ the corresponding offset operator (multiple channels form
their own Cartesian product). Power levels scale the physical waveform,
$c=p_i u$, so by the chain rule the reported gradient acquires a factor
$p_i$ and the Hessian $p_i^2$. Phase cycling multiplies the initial
state, the target state, and each X/Y control pair by case-specific
phases (Section~\ref{sec:phasecycle}), with the inverse phases applied to
the returned gradient and Hessian. Distortion chains and their Jacobians
are described in Section~\ref{sec:distortions}.

Three \emph{ensemble correlations} (\code{control.ens\_corrs}) restrict
the Cartesian product to physically correlated combinations:
\code{rho\_ens} gives each surviving case its own state pair,
\code{rho\_drift} pairs state $i$ with drift $i$, and \code{power\_drift}
pairs power $i$ with drift $i$. An \emph{ensemble budget}
(\code{control.budget}) draws a reproducible random subset of the full
case list --- either a member count or a fraction --- for stochastic
robustness at fixed cost. The bookkeeping lives in a six-column integer
\emph{catalog} built once at problem setup by \code{ens\_catalog.m}
(Section~\ref{sec:architecture}); the ensemble average runs over the
catalog rows, and summation order is fixed, so results are deterministic
and independent of the parallel pool size.

\subsection{Penalty functionals}
\label{sec:penalties}

Instrumental limits are enforced softly, by subtracting weighted penalty
terms from the fidelity, Eq.~\eqref{eq:objective}. With the normalised
waveform $u_{kn}$ ($K$ channels, $N_t$ columns), the implemented
functionals are:
\begin{align}
  P_{\mathrm{NS}}&=\frac{1}{N_t}\sum_{k,n}u_{kn}^{2}
  &&\text{norm square: favours low power;}\\
  P_{\mathrm{DNS}}&=\frac{1}{N_t}\sum_{k,n}
    \left[(\mathbf{u}\mathbf{D}^{\mathrm{T}})_{kn}\right]^{2}
  &&\text{derivative norm square: favours smoothness,}
\end{align}
where $\mathbf{D}$ is a five-point second-derivative matrix with wall
boundary conditions;
\begin{equation}
  P_{\mathrm{SNS}}=\frac{1}{N_t}\sum_{k,n}
  \left[(u_{kn}-b^{\mathrm{ceil}}_{kn})^2\,\theta(u_{kn}>b^{\mathrm{ceil}}_{kn})
       +(u_{kn}-b^{\mathrm{floor}}_{kn})^2\,\theta(u_{kn}<b^{\mathrm{floor}}_{kn})
  \right]
\end{equation}
is the \emph{spillout norm square}: zero inside the bounds
$[\,b^{\mathrm{floor}},b^{\mathrm{ceil}}\,]$ (scalars or per-element
arrays, in fractions of the power level) and quadratic outside --- the
default penalty, with weight 100, keeping the waveform inside the
amplifier range without biasing it inside. $P_{\mathrm{SNSA}}$ is SNS
applied to the \emph{amplitude} $\sqrt{u_X^2+u_Y^2}$ of each X/Y channel
pair after a polar transformation (with the exact polar chain rules for
gradient and Hessian): it caps the pulse amplitude while leaving the
phase free, which is the physically right constraint for
amplitude-limited, phase-agile hardware. All penalties return analytic
gradients and Hessians, which the optimiser combines with the fidelity
derivatives; penalties act on the ideal normalised waveform, before power
scaling and before any distortion chain.

\subsection{Stroboscopic steady state}
\label{sec:steady}

For experiments that repeat a pulse--and--delay block indefinitely (e.g.\
steady-state DNP), the object of optimisation is not a single transfer
but the \emph{stroboscopic steady state} of the period map. With the
period propagator $\mathbf{P}_{\mathrm{tot}}=\mathbf{P}_N\cdots
\mathbf{P}_1$ (keyholes excluded --- the map must be physical), the
steady state $\boldsymbol{\rho}_{\mathrm{ss}}$ satisfies
$\mathbf{P}_{\mathrm{tot}}\boldsymbol{\rho}_{\mathrm{ss}}
=\boldsymbol{\rho}_{\mathrm{ss}}$ and is found by the kernel
\code{steady()} routine; relaxation must be present (in Liouville space
with a thermal fixed point) for the map to be contractive. The fidelity
is the overlap of a zero-trace observable $\boldsymbol{\sigma}$ with the
steady state at the stroboscopic grid, and its gradient requires the
derivative of the fixed point. Writing the steady state as the geometric
series of period maps acting on the unit-trace component shows that the
correct adjoint seed is the \emph{dressed target}
\begin{equation}
  \boldsymbol{\chi}_0:\quad
  (\mathbf{1}-\mathbf{P}_{\mathrm{tot}})^{\dagger}\,\boldsymbol{\chi}_0
  =\boldsymbol{\sigma}
  \quad\text{on the traceless subspace,}
\end{equation}
solved with the unit-state singularity excluded (the code works on the
subspace orthogonal to the identity; the target must therefore be
traceless, which is checked). With $\boldsymbol{\rho}_0$ overwritten by
$\boldsymbol{\rho}_{\mathrm{ss}}$ and $\boldsymbol{\sigma}$ by
$\boldsymbol{\chi}_0$, the ordinary GRAPE gradient machinery of
Section~\ref{sec:grapeliouv} then delivers the steady-state fidelity
derivative. Steady-state mode requires the \code{sphten-liouv} formalism,
the rectangle integrator, first-order optimisers, and no
prefix/suffix/dead-time decorations; the initial state supplied by the
user is ignored (with a warning), because the steady state does not
depend on it.

\subsection{Phase cycles and cooperative pulses}
\label{sec:phasecycle}

Phase cycling is built into the ensemble machinery: each line of
\code{control.phase\_cycle} is
$[\varphi_{\mathrm{init}},\ \varphi_1,\ldots,\varphi_{K/2},\
\varphi_{\mathrm{targ}}]$, and the corresponding ensemble case runs with
$\boldsymbol{\rho}_0\to e^{i\varphi_{\mathrm{init}}}\boldsymbol{\rho}_0$,
$\boldsymbol{\sigma}\to e^{i\varphi_{\mathrm{targ}}}\boldsymbol{\sigma}$,
and each X/Y control pair rotated through its phase,
$(u_X+iu_Y)\to e^{i\varphi_k}(u_X+iu_Y)$. The returned gradient (and
Hessian) is rotated back through $e^{-i\varphi_k}$, which is the exact
chain rule of the rotation. Maximising the phase-cycled average fidelity
produces pulses whose \emph{coherence-pathway-selected sum} behaves as
specified.

Cooperative pulses~(\code{grape\_coop}) push this one step further for
phase-modulated point-to-point transfers: two pulses are optimised
simultaneously so that each maximises its own fidelity while their
\emph{impurities} --- the projections of the final states off the target
--- cancel in the sum of the two experiments. The functional is
\begin{equation}
  F=\frac{f_a+f_b}{2}
  -\left\langle\left\|
   (\mathbf{1}-\mathbf{P}_{\sigma})
   (\boldsymbol{\rho}_a(T)+\boldsymbol{\rho}_b(T))
   \right\|^2\right\rangle,
\end{equation}
with $\mathbf{P}_{\sigma}$ the projector onto the target state and the
average running over the ensemble; the impurity gradient is assembled by
a second pass through the phase-GRAPE machinery with the summed impurity
as the target. The two phase profiles travel side by side in one
optimisation vector.
%===============================================================
\section{Numerical implementation in Spinach}
%===============================================================

\subsection{Propagation primitives}

All optimal control code is built on the two kernel propagation
primitives. Given a generator and an interval, \code{propagator()}
returns the matrix $\exp(-i\mathbf{L}\Delta t)$ using
scaled-and-squared Taylor exponentiation with aggressive sparsity
bookkeeping; small elements are chopped at \code{tols.prop\_chop}, and
computed propagators are cached (in memory and, for large objects, on
disk) so that repeated requests are cheap. \code{step()} computes the
\emph{action} $\exp(-i\mathbf{L}\Delta t)\,\boldsymbol{\rho}$ on a
state without ever forming the propagator, which is the only practical
option at Liouville-space dimensions; given a cell array of edge and
midpoint generators it applies the piecewise-linear product quadrature
of Section~\ref{sec:trapezium}. The optimal control module uses
\code{propagator()} where matrices are genuinely needed (Hilbert space,
steady-state period maps, Goodwin cumulative propagators) and
\code{step()} everywhere else.

The auxiliary matrices of Section~\ref{sec:dirdiff} are assembled as
sparse blocks --- the doubled or tripled dimension is harmless because
the top-right coupling block is as sparse as the control operator ---
and either exponentiated with \code{propagator()} (Hilbert space, via
\code{dirdiff.m}) or applied to stacked vectors with \code{step()}
(Liouville space). Propagation inside the derivative loops runs with
reporting hushed; the outer functions restore the console before
printing.

Numerical hygiene points that matter in practice:
\begin{itemize}
\item Control operators, offset operators, and drifts are cleaned up
  (\code{clean\_up}) at problem setup, so numerical dust in the inputs
  does not degrade the sparsity of the auxiliary matrices.
\item Where unitary propagation is mathematically required --- Hilbert
  space, wavefunction formalism, and the Goodwin Hessian --- all
  generators are checked for Hermiticity at setup and refused otherwise.
\item Exactly zero fidelity or an exactly zero gradient aborts the run
  with an explanatory error: the target is unreachable from the source
  (a symmetry is in the way), or the initial guess is pathologically
  poor. This check has saved a great many silently useless optimisations.
\item The first LBFGS iteration takes a deliberately conservative step,
  $0.01\,\mathbf{g}/\max|\mathbf{g}|$, because the fidelity landscape
  near a random guess is often violently nonlinear.
\end{itemize}

\subsection{The optimiser stack}
\label{sec:optimisers}

The maximiser \code{fmaxnewton(spin\_system,cost\_function,guess)} runs
the iteration loop for all methods; the objective is called through
\code{objeval()}, which assembles the total functional of
Eq.~\eqref{eq:objective} from the fidelity-and-penalties vector returned
by the cost function, keeps evaluation counters, and reshapes between
waveform arrays and optimisation vectors.

\begin{description}
\item[\code{lbfgs}] (default): limited-memory BFGS ascent
  \cite{liu1989}. The inverse Hessian is applied implicitly from a
  history of waveform and gradient differences; the history depth is
  \code{control.n\_grads} (default 50). No regularisation is possible or
  needed; memory cost is linear in problem size.
\item[\code{rbfgs}]: full-matrix BFGS pseudo-Hessian, regularised before
  inversion (below). Useful when the pseudo-Hessian conditioning, not the
  memory, is the limiting factor.
\item[\code{newton}]: analytic Hessian of
  Section~\ref{sec:hessians}, symmetrised and regularised. Quadratic
  local convergence; each iteration costs one Hessian evaluation.
\item[\code{goodwin}]: the same Newton step with the cross-interval
  Hessian assembled through cumulative propagators
  \cite{goodwin2016}; requires Hermitian generators.
\end{description}

Newton-family Hessians are made safely negative-definite (this is a
maximisation) by \emph{rational function optimisation} regularisation
\cite{banerjee1985} in the kernel routine \code{hessreg()}: the Hessian
is embedded in an augmented matrix whose scaled eigenvalue problem shifts
the spectrum away from zero, with the scaling iterated (up to
\code{reg\_max\_iter}, condition number capped at \code{reg\_max\_cond})
until the step is well-posed. If a computed direction still fails the
ascent test $\mathbf{g}^{\mathrm{T}}\mathbf{d}>0$, the optimiser falls
back to projected steepest ascent for that iteration.

The step length is found by a Fletcher-style two-stage line search
\cite{fletcher2000}: \code{bracketing()} expands an interval that is
guaranteed to contain acceptable points, and \code{sectioning()}
contracts it with cubic interpolation until the strong Wolfe conditions
hold --- sufficient increase with $c_1=10^{-2}$ and curvature with
$c_2=0.9$ (bracket expansion factor 3, section contraction bounds 0.1 and
0.5; these live in \code{control.ls\_*} and rarely need attention).
Termination is on step norm ($\|\alpha\mathbf{d}\|_1<$ \code{tol\_x},
default $10^{-3}$), gradient norm ($\|\mathbf{g}\|_2<$ \code{tol\_g},
default $10^{-6}$), iteration count (\code{max\_iter}, default 100), or a
failed line search. Every iteration prints fidelity, penalties, total,
step length, gradient norm, and the angle between the search direction
and the gradient (SDGA) --- a useful health indicator: it should be well
below $90^\circ$ and typically shrinks as curvature information
accumulates.

Optional machinery hanging off the loop: a checkpoint file
(\code{control.checkpoint}) saves the current waveform in the scratch
directory after every accepted step and restarts from it if the run is
interrupted; a video writer (\code{control.video\_file}) grabs the
diagnostic figure every iteration; a freeze mask
(\code{control.freeze}) excludes selected waveform elements from the
optimisation everywhere --- gradient, Hessian, history, and line search.

\subsection{Coordinate systems and wrappers}
\label{sec:wrappers}

The kernel GRAPE functions differentiate with respect to raw control
coefficients; user-facing parametrisations are handled by thin wrappers
that apply exact chain rules:

\begin{description}
\item[\code{grape\_xy}] --- Cartesian X/Y control amplitudes, the
  standard entry point, called as the cost function of
  \code{fmaxnewton}. If a waveform basis $\mathbf{B}$ is supplied
  (\code{control.basis}, orthonormal rows), the optimisation variables
  are the expansion coefficients $\mathbf{w}$, the physical waveform is
  $\mathbf{u}=\mathbf{w}\mathbf{B}$, and gradients pull back through
  $\mathbf{B}^{\mathrm{T}}$; Hessians are reordered and transformed with
  the corresponding tensor contraction. Penalties are computed here and
  returned in the separated fidelity vector.
\item[\code{grape\_phase}] --- phase-modulated pulses with a fixed
  amplitude profile (\code{control.amplitudes}): variables are the
  phases $\varphi_k(t)$, the Cartesian waveform is
  $(A\cos\varphi,A\sin\varphi)$, and the polar chain rule converts
  gradients and Hessians; only the phase gradient is returned.
\item[\code{grape\_curv}] --- arbitrary user-defined curvilinear
  coordinates: the user supplies $u\!\to\!x$ maps and their Jacobians,
  and the wrapper applies them; penalties are evaluated in the
  rectilinear representation.
\item[\code{grape\_coop}] --- cooperative pulse pairs
  (Section~\ref{sec:phasecycle}) on top of \code{grape\_phase}.
\end{description}

A separate special-purpose objective, \code{tgrape()}, differentiates the
fidelity with respect to the \emph{interval durations} at a fixed
waveform: since
$\partial\exp(-i\mathbf{L}\Delta t)/\partial\Delta t
=-i\mathbf{L}\exp(-i\mathbf{L}\Delta t)$, the duration gradient needs
only generator actions on the stored trajectory. Durations are passed in
a user-chosen time unit so that the optimisation variables are of order
one --- optimisers stall on badly scaled variables.

\subsection{Diagnostics and plotting}
\label{sec:plotting}

If \code{control.plotting} is non-empty, the client calls
\code{ctrl\_trajan()} after every objective evaluation with the current
waveform, trajectory data, and per-case fidelities. Available panels:
\begin{center}
\begin{tabular}{lp{8.6cm}}
\toprule
waveform & \code{xy\_controls}, \code{phi\_controls},
           \code{amp\_controls}, \code{frq\_controls}
           (instantaneous frequency)\\
trajectory & \code{trajectory}, \code{correlation\_order},
             \code{coherence\_order}, \code{local\_each\_spin},
             \code{total\_each\_spin}, \code{level\_populations}\\
ensemble & \code{robustness} (fidelity across the ensemble)\\
time-frequency & \code{spectrogram}, \code{time\_by\_slice}\\
\bottomrule
\end{tabular}
\end{center} Trajectory-derived panels force the (large) trajectory stack to
be returned from the workers, and all plotting runs on the client every
single evaluation --- hence the setup-time warning that diagnostic plots
are computationally expensive. \code{control.distplot} optionally applies
a distortion chain to the waveform before plotting, so that the display
shows what the hardware would emit.

\subsection{Waveform distortions}
\label{sec:distortions}

Real pulses are distorted between the waveform generator and the sample
--- amplifier compression, resonator ringdown, transmission line
response. Spinach handles this exactly: \code{control.distortion} is a
cell array of function handles, each mapping a physical waveform to a
distorted physical waveform and, on request, returning the Jacobian of
that map. During every ensemble case the chain is applied in sequence,
\begin{equation}
  \mathbf{u}_{\mathrm{phys}}
  =D_M(\cdots D_2(D_1(p\,\mathbf{u}))\cdots),
  \qquad
  \mathbf{J}=\mathbf{J}_M\cdots\mathbf{J}_2\mathbf{J}_1,
\end{equation}
and the gradient returned by GRAPE is pulled back through
$\mathbf{J}^{\mathrm{T}}$ before the power-level chain rule. Rows of the
cell array form an ensemble dimension (distinct distortion realisations
to be robust against); columns form the chain stages. Distortion
functions may implement anything differentiable --- RLC response models,
amplifier saturation curves, measured impulse responses --- and the
supplied Jacobian is verified by the derivative test examples. Because
the map composes with the fidelity, only first-order optimisers are
permitted when distortions are present (\code{lbfgs}); analytic Hessians
of composed distortion chains are not implemented.

%===============================================================
\section{Spinach optimal control module architecture}
\label{sec:architecture}
%===============================================================

\subsection{Data flow}

An optimal control calculation passes through the following layers, top
to bottom:
\begin{center}
\begin{tabular}{ll}
\toprule
layer & role \\
\midrule
\code{optimcon()} & parses and freezes the problem, builds the catalog,\\
                  & publishes worker-resident invariants\\
\code{fmaxnewton()} & optimiser iterations and line search\\
\code{objeval()} & objective assembly, counters, reshaping\\
\code{grape\_xy()} and other wrappers
                  & coordinate chain rules, penalties\\
\code{ensemble()} & parallel loop over ensemble cases\\
\code{grape\_liouv()}, \code{grape\_hilb()} & trajectories and
                    propagator derivatives\\
\code{step()}, \code{propagator()}, & \\
\code{dirdiff()}, \code{aux\_mat()} & propagation primitives\\
\bottomrule
\end{tabular}
\end{center}
The user builds the \code{control} structure, calls
\code{spin\_system=optimcon(spin\_system,control)} once, and hands
\code{fmaxnewton} a cost function --- almost always \code{@grape\_xy} or
\code{@grape\_phase} --- together with an initial guess.

\subsection{The frozen problem and the worker-resident ensemble}

The architectural principle of the module is:
\emph{the ensemble is the only parallel dimension, and the
evaluation-invariant problem data crosses the process boundary exactly
once per optimisation.}

\code{optimcon()} ends by freezing the problem. It builds the
\emph{ensemble case catalog} --- the six-column integer array of
Section~\ref{sec:ensembles}, one row per case, produced by
\code{ens\_catalog()} together with the dimension sizes --- and then
publishes the complete frozen problem (the entire \code{spin\_system},
control structure included) to the parallel pool as a
\code{parallel.pool.Constant} held in \code{control.invariants}. The
heavy invariant fields --- drifts, control operators, offset operators,
the trapezium control--control commutator tables, and the Bloch--Siegert
response operators, when present --- are then removed from the
client-side structure, and their names are recorded in
\code{control.frozen\_fields}. A
\code{parallel.pool.Constant} constructs without a pool, broadcasts
lazily to pools created later, and survives pool deletion and recreation,
so serial runs and delayed pool startup work without special cases.

Each objective evaluation calls \code{ensemble()}, which runs a
\code{parfor} over the catalog rows. Every case fetches the
worker-resident frozen problem (a fast local dereference), replaces its
control structure wholesale with the live client-side one, re-attaches
the fields named in \code{control.frozen\_fields}, assembles its
drift (offset terms added per the catalog), applies its phase cycle and
distortion chain, calls the serial GRAPE kernel, and returns fidelity,
gradient, and Hessian contributions. The client sums the contributions in
fixed order and divides by the case count. Per-case propagation is pure
serial: the eleven interior \code{parfor} loops of the previous
architecture are gone, along with the time-parallel strategy and all
\code{ValueStore} traffic.

The consequences, measured on a Liouville dimension 21067 problem with
36 ensemble cases and 12 process workers: interprocess traffic fell from
96.97 to 0.20~MB per objective evaluation, and wall clock from 79.9 to
60.8~s median --- the per-evaluation broadcast serialisation was costing
about a quarter of the wall time at scale. Only the waveform, the states,
and small numerics travel at each evaluation.

Because cases are distributed over workers whole, the parallel
efficiency ceiling is the load balance
\begin{equation}
  \eta_{\max}
  =\frac{M}{W\left\lceil M/W\right\rceil},
\end{equation}
for $M$ cases on $W$ workers; \code{optimcon()} prints this number at
setup. It pays to make the case count a multiple of the worker count
(e.g.\ 36 cases on 12 workers gives $\eta_{\max}=1$; 37 cases give
$0.77$). Single-case optimisations do not parallelise --- for those, the
serial kernel is already the optimum at this granularity.

\subsection{The mutation contract}
\label{sec:contract}

Freezing the problem draws a sharp line through the
\code{spin\_system.control} structure that every user of sweep drivers
should know:

\begin{description}
\item[May be overwritten between optimisations] (the \emph{live} set):
  every control field not named in \code{control.frozen\_fields} ---
  in particular \code{rho\_init}, \code{rho\_targ} (same member
  counts), \code{pulse\_dt} (same length), \code{pwr\_levels},
  \code{offsets} (same value counts), \code{phase\_cycle} (same line
  count), \code{distortion} (same rows), \code{freeze},
  \code{fidelity}, \code{plotting}, \code{keyholes}, penalties,
  weights, bounds, and the prefix and suffix functions. The live
  structure travels to the workers fresh at every objective
  evaluation, so overwrites take effect at the next one --- this is
  how parameter sweeps reuse one \code{optimcon()} setup across many
  optimisations.
\item[Frozen] (changes require a fresh \code{optimcon()} call): the
  ensemble \emph{composition} --- the member counts of every ensemble
  dimension, the correlation settings \code{ens\_corrs}, and the budget
  --- and the worker-resident invariants named in
  \code{control.frozen\_fields}: \code{drifts}, \code{operators},
  \code{off\_ops}, the trapezium commutator tables \code{cc\_comm} and
  \code{cc\_comm\_idx}, and the Bloch--Siegert \code{resp\_ops}. These
  fields are physically absent from the
  client structure after \code{optimcon()}; touching them, or changing
  anything that alters the case catalog, raises an immediate error
  rather than silently running a stale ensemble. The guard literally
  recomputes the catalog from the current control structure at every
  evaluation and requires it to be identical to the frozen one, so any
  catalog-defining drift is caught, including settings added in the
  future.
\end{description}

Two further architectural guarantees. First, \emph{callbacks see
everything}: prefix and suffix functions execute on the workers against
the complete frozen problem with the live fields grafted fresh, so a
callback that reads \code{control.parameters} --- the designated stash
for user data, accepted by \code{optimcon()} without parsing --- or any
other control field behaves exactly as it would client-side. Second,
\emph{direct low-level calls keep working}: \code{grape\_liouv()} and
\code{grape\_hilb()} accept the client-side structure with explicit
drift and control arguments, which is how the shipped regression tests
drive them.

\subsection{Parallel execution model}

The module never opens or resizes pools on its own: it uses the pool
that exists (\code{sys.parallel} at \code{create()} time, or an explicit
\code{parpool} call), and runs serially if there is none ---
\code{parfor} with a zero worker cap degrades gracefully. Determinism is
worker-count-independent because summation happens client-side in
catalog order. The \code{parallel.pool.Constant} transport is documented
by MathWorks for process-local and cluster pools alike; local process
pools are the measured configuration.

For very large ensembles, the budget mechanism
(Section~\ref{sec:ensembles}) bounds the per-evaluation cost; the random
subset is drawn once, at setup, with a fixed seed, so the optimisation
surface is deterministic. For very large state spaces with a single
case, parallelism must come from within the propagation primitives
(threaded BLAS, GPU arrays supplied as drifts and controls), not from
the ensemble loop.
%===============================================================
\section{Getting started with optimal control in Spinach}
%===============================================================

\subsection{A complete first example}

The following script designs a pulse that moves longitudinal
magnetisation from the proton to the carbon in a two-spin system through
the $J$-coupling, using controls on both channels --- the classic
state-to-state transfer. It is complete and runnable as shown.

\begin{lstlisting}
% Parallel pool: eight local process workers
sys.parallel={'processes',8};

% Magnet and spin system
sys.magnet=9.4;
sys.isotopes={'1H','13C'};

% Chemical shifts, ppm, and J-coupling, Hz
inter.zeeman.scalar={1.0,50.0};
inter.coupling.scalar=cell(2);
inter.coupling.scalar{1,2}=140;

% Basis set: full spherical tensor basis, Liouville space
bas.formalism='sphten-liouv';
bas.approximation='none';

% Spinach housekeeping
spin_system=create(sys,inter);
spin_system=basis(spin_system,bas);

% Drift Liouvillian under the NMR rotating-frame assumptions
H=hamiltonian(assume(spin_system,'nmr'));

% Normalised initial and target states
rho_init=state(spin_system,{'Lz'},{1});
rho_init=rho_init/norm(full(rho_init),2);
rho_targ=state(spin_system,{'Lz'},{2});
rho_targ=rho_targ/norm(full(rho_targ),2);

% Control operators: X and Y on both channels
LxH=operator(spin_system,'Lx','1H');
LyH=operator(spin_system,'Ly','1H');
LxC=operator(spin_system,'Lx','13C');
LyC=operator(spin_system,'Ly','13C');

% Control problem specification
control.isotopes={'1H','13C'};                 % isotopes under control
control.channels=[1;1;2;2];                    % channel of each operator
control.drifts={{H}};                          % drift ensemble: one member
control.operators={LxH,LyH,LxC,LyC};           % four control channels
control.rho_init={rho_init};                   % one state pair
control.rho_targ={rho_targ};
control.pwr_levels=2*pi*linspace(9e3,11e3,3);  % B1 ensemble, rad/s
control.pulse_dt=1e-4*ones(1,50);              % 50 steps of 100 us
control.method='lbfgs';                        % optimiser
control.max_iter=100;                          % iteration cap
control.plotting={'xy_controls','robustness'}; % live diagnostics

% Freeze the problem
spin_system=optimcon(spin_system,control);

% Random normalised initial guess
guess=randn(4,50)/10;

% Run the optimisation
pulse=fmaxnewton(spin_system,@grape_xy,guess);
\end{lstlisting}

Everything before the \code{control} block is ordinary Spinach setup; the
optimal control part is the control structure, one \code{optimcon()}
call, and one \code{fmaxnewton()} call. The result \code{pulse} is the
optimised normalised waveform, one row per control operator, in the same
layout as the guess.

\subsection{Anatomy of the control structure}

Eight fields are mandatory:
\begin{center}
\begin{tabular}{lp{9.5cm}}
\toprule
field & content \\
\midrule
\code{drifts} & drift generators: a cell array over the ensemble of cell
arrays over time --- \code{\{\{H\}\}} for one static drift;
\code{\{\{H1\},\{H2\},...\}} for a drift ensemble; each inner array has
either one element or one per time point for time-dependent drifts\\
\code{operators} & cell array of control operators\\
\code{isotopes} & cell array of isotope strings: the isotopes that the
control channels address\\
\code{channels} & vector with one element per control operator: the
index into \code{isotopes} of the channel that the operator drives\\
\code{rho\_init} & cell array of initial states (vectors in Liouville
space, density matrices in Hilbert space)\\
\code{rho\_targ} & cell array of target states, same length\\
\code{pwr\_levels} & row vector of nominal control powers, rad/s ---
more than one value creates a $B_1$ robustness ensemble\\
\code{pulse\_dt} & row vector of interval durations, seconds\\
\bottomrule
\end{tabular}
\end{center}
Everything else has a sensible default (Section~\ref{sec:reference}).
The waveform the optimiser works with is normalised: a value of $1$ on a
channel means the full nominal power. The guess should therefore be
order-of-tenths random numbers or a physically motivated shape scaled the
same way; the number of rows equals the number of control operators, the
number of columns equals the number of intervals
(\code{numel(pulse\_dt)}) for the rectangle integrator, one more for the
trapezium integrator.

\subsection{Reading the console output}

\code{optimcon()} prints one line per parsed option, the ensemble case
count, the pool size, and the parallel efficiency ceiling ---
worth a glance before a long run. \code{fmaxnewton()} then prints one
line per iteration:
\begin{center}
\begin{tabular}{ll}
\toprule
column & meaning \\
\midrule
\code{Iter} & iteration number\\
\code{\#f, \#g, \#H, \#R} & cumulative objective, gradient, Hessian,
                            and regularisation counts\\
\code{fidelity} & current fidelity $f$ \\
\code{penalties} & current weighted penalty total\\
\code{total} & the maximised objective, $f-\sum w_jP_j$\\
\code{alpha} & accepted line search step\\
\code{|grad|} & gradient 2-norm\\
\code{SDGA} & angle (degrees) between search direction and gradient\\
\bottomrule
\end{tabular}
\end{center}
A healthy run shows the total climbing, the gradient norm falling, and
SDGA well below $90^\circ$. The footer reports which termination
condition fired. For a state-to-state transfer with the default
\code{real} fidelity, $f=1$ is a perfect transfer; how close one can get
depends on the pulse duration relative to the coupling topology and on
the ensemble spread.

\subsection{Using the result}

The returned waveform is normalised; the physical pulse on channel $k$
is \code{pwr\_levels(m)*pulse(k,:)} for the power level of interest.
From here:
\begin{itemize}
\item simulate the pulse in any Spinach sequence by adding
  $\sum_k c_k(t)\,\mathbf{C}_k$ to the drift interval by interval, or by
  calling the shaped-pulse machinery of the kernel;
\item export amplitude and phase per channel:
  $A=\sqrt{u_X^2+u_Y^2}\cdot p/(2\pi)$ in Hz and
  $\varphi=\operatorname{atan2}(u_Y,u_X)$, then resample to the
  instrument grid;
\item verify robustness by evaluating the fidelity over a denser
  offset/power grid than was optimised over (rebuild the control
  structure with the dense ensemble, rerun \code{optimcon()}, and call
  the cost function once with \code{max\_iter=0} or evaluate
  \code{grape\_xy} directly);
\item continue the optimisation from the result --- the returned
  waveform is a valid guess, including across changes of live settings
  (Section~\ref{sec:contract}).
\end{itemize}

%===============================================================
\section{Advanced user manual}
%===============================================================

\subsection{Robustness ensembles}

The six ensemble dimensions compose freely; the case count is their
product (after correlations and budget). The two workhorse dimensions:

\emph{$B_1$ miscalibration.} Several entries in \code{pwr\_levels}
optimise the average fidelity over amplifier miscalibration; three to
five values spanning $\pm10\%$ around nominal are typical.

\emph{Resonance offset.} \code{control.offsets} is a cell array of row
vectors of offset values in Hz, one cell per offset channel, with the
matching operators in \code{control.off\_ops} (typically $2\pi$-free
$\op{L}_Z$ operators of the irradiated species --- the code multiplies by
$2\pi$ internally). Multiple channels form a Cartesian product. For a
broadband proton pulse: \code{control.offsets=\{linspace(-5000,5000,25)\}},
\code{control.off\_ops=\{LzH\}}.

For drift ensembles (orientations of a powder, conformer sets, hardware
$B_0$ maps), supply one drift per member; the \code{rho\_drift} and
\code{power\_drift} correlations bind the state or power dimension to the
drift dimension when members are physically linked (e.g.\ each crystal
orientation carrying its own starting state). When the full product is
too expensive, \code{control.budget} caps the evaluated subset: an
integer is a member count, a value $\leq1$ is a fraction; the subset is
drawn once with a fixed seed, so the objective stays deterministic.
Optimising on a budget subset and \emph{validating on the full grid} is
standard practice.

\subsection{Time-dependent drifts}

Each drift ensemble member may be a cell array with one generator per
time point, which is how magic angle spinning, field sweeps, and other
programmed drift modulations enter: the propagation cycles through the
inner array as it steps through the pulse. All members must have the same
number of time slices, and that number must equal the control array
length. Rotor-synchronised optimal control under MAS is set up exactly
this way.

\subsection{Phase cycles}

\code{control.phase\_cycle} requires an even number of controls arranged
in X/Y pairs. Each row is one cycle line:
$[\varphi_{\mathrm{init}},\varphi_1,\ldots,\varphi_{K/2},
\varphi_{\mathrm{targ}}]$, radians. The optimiser then maximises the
average fidelity of the cycle, i.e.\ designs the pulse so that the
\emph{sum} of the phase-shifted experiments performs the transfer while
the unwanted pathways cancel. Note that the number of columns is
$K/2+2$; a receiver phase is folded into $\varphi_{\mathrm{targ}}$.

\subsection{Waveform distortions}

Supply \code{control.distortion} as an $E\times S$ cell array of function
handles: $S$ chain stages applied left to right, $E$ ensemble rows. Each
handle must implement
\begin{lstlisting}
[wf_out,J]=my_distortion(wf_in)   % J optional, needed for gradients
\end{lstlisting}
on the \emph{physical} (power-scaled) waveform, returning the Jacobian of
the stacked waveform vector on request. Shipped examples cover RLC
resonator response chains and the associated Jacobians; the derivative
test suite (\code{examples/fundamentals/derivative\_tests}) shows how to
verify a hand-written Jacobian against finite differences before trusting
an optimisation to it. Distortions restrict the optimiser to
\code{lbfgs}, and combine freely with all other ensemble dimensions ---
distinct rows model unit-to-unit hardware spread.

\subsection{Keyholes, prefix, suffix, and dead time}

\emph{Keyholes} (\code{control.keyholes}, rectangle integrator only) are
projectors applied to the state after every interval: a cell array with
one entry per time point, each empty or a linear idempotent function
handle (both properties are verified numerically at setup). They
implement engineered subspace constraints --- decoherence-protected
manifolds, coherence-order gating --- directly in the trajectory and its
derivatives.

\emph{Prefix and suffix} (\code{control.prefix} and
\code{control.suffix}) are function handles with the signature
\begin{center}
\code{rho=f(spin\_system,drift,rho)}
\end{center}
applied to the initial state before the pulse and to the target state
before the backward propagation (coded in reverse time), used for fixed
preparation and readout blocks around the optimised segment. They execute on the parallel workers
against the complete control structure, so they may read any control
field, including the free-form \code{control.parameters} stash.

\emph{Dead time} (\code{control.dead\_time}) evolves the target backward
under the last drift for a fixed delay, producing pulses whose outcome
survives a receiver dead time; it is applied with the adjoint generator,
so relaxation during the dead time is handled correctly.

\subsection{Trapezium integrator in practice}

Set \code{control.integrator='trapezium'} and give the waveform (and
guess) $N+1$ columns for $N$ intervals. The optimised pulse is then
natively piecewise-linear --- what the instrument plays between DAC
points --- and the fidelity model error drops from second to fourth
order in the interval length~\cite{rasulov2023}, which allows far
coarser time grids at equal accuracy. Restrictions: first-order
optimisers only (no analytic Hessian), no keyholes, and no steady-state
mode. Control--control commutators are precomputed at setup; for
commuting control sets (single-channel X/Y pairs commute with themselves
only trivially) the extra cost over the rectangle integrator is small.

\subsection{Stroboscopic steady-state optimisation}

Set \code{control.steady=true()}. Requirements: \code{sphten-liouv}
formalism, relaxation present, rectangle integrator, first-order
optimiser, traceless target, no prefix/suffix/dead time; the supplied
initial state is ignored. The optimised functional is the overlap of the
target observable with the periodic steady state of the
pulse--and--relaxation cycle --- the natural objective for steady-state
DNP and related experiments; the \code{steady\_orbit} example family
shows complete solid effect setups.

\subsection{Bloch--Siegert corrections}
\label{sec:bsiegert}

A resonant control operator is the rotating-frame remnant of a linearly
polarised laboratory-frame drive. The counter-rotating component that
the rotating-wave approximation discards shifts every energy level in
second order of the drive amplitude, and the same drive also shifts
every spin that it is not resonant with. Setting
\code{control.bsiegert=true()} keeps these shifts in the optimisation:
during problem freezing, \code{optimcon()} calls \code{bss\_ops()} to
build one response operator per control channel,
\begin{equation}
  \mathbf{B}_k=\sum_{n\,\in\,k}
  \frac{\mathbf{L}_{\mathrm{Z}}^{(n)}}{2\left(\omega_n+\omega_k\right)}
  +\sum_{n\,\notin\,k}
  \left(\frac{\gamma_n}{\gamma_k}\right)^{2}
  \frac{\omega_n\,\mathbf{L}_{\mathrm{Z}}^{(n)}}
       {\omega_n^{2}-\omega_k^{2}}
\end{equation}
where $n\in k$ runs over the spins of the isotope that the $k$-th
channel addresses, $\omega_n$ are the signed laboratory-frame Zeeman
frequencies, and $\omega_k$ and $\gamma_k$ are the carrier frequency
and the magnetogyric ratio of the channel isotope. Spins of the channel
isotope receive only the never-resonant term, because their resonant
term is the control operator itself, which GRAPE propagates exactly;
for every other spin the resonant and the never-resonant terms sum into
the familiar off-resonant shift
$\omega_n\omega_{1n}^{2}/(\omega_n^{2}-\omega_k^{2})$, in which
$\omega_{1n}=(\gamma_n/\gamma_k)c_{kn}$ is the nutation frequency the
drive induces at spin $n$. The slice generators acquire a quadratic
term,
\begin{equation}
  \mathbf{L}_n=\mathbf{L}_0
  +\sum_k c_{kn}\mathbf{C}_k
  +\sum_k c_{kn}^{2}\,\mathbf{B}_k
\end{equation}
and the derivative direction $\mathbf{C}_k+2c_{kn}\mathbf{B}_k$
replaces $\mathbf{C}_k$ in the gradient calculation --- the gradient
stays analytically exact.

The waveform coefficients $c_{kn}$ must therefore be physical nutation
frequencies: \code{optimcon()} refuses control operators on corrected
channels that are not unit quadratures
$\cos\varphi\,\mathbf{L}_{\mathrm{X}}+\sin\varphi\,\mathbf{L}_{\mathrm{Y}}$
of their channel isotope. Carrier frequencies are taken on resonance
with each channel isotope at the current magnet field; zero and
degenerate carrier denominators are refused, as are the trapezium
integrator and the exact-Hessian optimisers. The response operators are
stored without clean-up because their matrix elements scale as inverse
Zeeman frequencies.

To replay an optimised pulse outside the optimiser under the same slice
generators, \code{bloch\_siegert()} appends the response operators to
the control operator list as virtual channels with squared coefficient
vectors, ready for the piecewise-constant methods of
\code{shaped\_pulse\_xy()}:

\begin{lstlisting}
% Amplitudes in rad/s from the normalised waveform
coefs=mat2cell(pwr*pulse,[1 1]);

% Append the Bloch-Siegert virtual channels
[ops,coefs]=bloch_siegert(spin_system,{Lx,Ly},coefs);

% Propagate with a piecewise-constant method
rho=shaped_pulse_xy(spin_system,H,ops,coefs,...
                    control.pulse_dt,rho,'expv-pwc');
\end{lstlisting}

The \code{examples/optimal\_control/bloch\_siegert} folder demonstrates
the corrections on transfer fidelities as a function of drive power, on
spectator-spin Ramsey phases, and on the classic 1940 level-shift
experiment; \code{difdiff\_bs\_rect.m} in the derivative test folder
verifies the corrected gradients against finite differences.

\subsection{Second-order methods}

\code{newton} and \code{goodwin} converge in far fewer iterations than
LBFGS once inside the quadratic basin, at a much higher per-iteration
cost ($O(K^2N)$ auxiliary propagations for the same-interval blocks plus
the cross-interval assembly, against $O(KN)$ for the gradient). They pay
off when the fidelity landscape is stiff --- strongly coupled targets,
tight bounds through penalties, ill-conditioned bases --- and on modest
state-space dimensions where propagator matrices are affordable. The RFO
regularisation makes raw saddle-ridden GRAPE Hessians usable; the
regularisation iteration count appears in the \code{\#R} column. Goodwin
assembly additionally requires all generators Hermitian and is the
faster of the two when its cumulative propagators fit in memory. A
practical pattern is LBFGS to $f\sim0.9$, then Newton to convergence ---
run one optimisation, overwrite \code{control.method} is \emph{not}
possible (method is frozen); instead re-run \code{optimcon()} with the
new method and pass the LBFGS result as the guess.

\subsection{Penalties, bounds, and freeze masks}

The default is SNS with weight 100 and bounds $\pm1$ --- soft clipping at
nominal power. To change the admissible corridor, set
\code{control.u\_bound}/\code{l\_bound} (scalars, or arrays shaped like
the waveform for time- and channel-dependent envelopes). For
phase-modulated pulses with a fixed amplitude profile, bounds are
meaningless (amplitude is fixed); for amplitude-limited phase-agile
hardware, use \code{SNSA}, which caps
$\sqrt{u_X^2+u_Y^2}$ instead of the Cartesian components. \code{DNS}
with a modest weight produces smooth, hardware-friendly waveforms;
combining \code{\{'SNS','DNS'\}} with weights like \code{[100 10]} is a
common choice. Weights are part of the objective: check the
\code{penalties} column stays small compared to the fidelity, otherwise
the optimiser is fighting the penalty rather than the physics.

\code{control.freeze} is a logical array of the waveform shape: true
elements are locked at their guess values --- used to pin pulse edges to
zero, protect hand-designed segments, or optimise one channel at a time.
For \code{grape\_phase}, the freeze mask has the phase-profile shape and
is expanded to the Cartesian pairs internally.

\subsection{Sweeps and re-optimisation}

The mutation contract (Section~\ref{sec:contract}) is designed for
sweep drivers: after one \code{optimcon()} call, loops may overwrite
live numerical fields --- pulse durations, powers, offset values, phase
cycles, states, penalties, bounds, plotting --- and re-run
\code{fmaxnewton()} each time, typically warm-starting from the previous
result. Composition changes (member counts, correlations, budget) and
generator changes error out immediately; rebuild the control structure
and call \code{optimcon()} again --- it is cheap relative to any
optimisation. For unattended long runs, set \code{control.checkpoint};
for post-mortem analysis of convergence, \code{control.video\_file}
records the diagnostic figure per iteration.

\subsection{Parallel performance tuning}

The rules of thumb that follow from the architecture:
\begin{itemize}
\item Aim for a case count that is a multiple of the worker count; the
  setup printout shows the achievable ceiling.
\item More workers than cases is waste; fewer, heavier cases per worker
  amortise the fixed per-evaluation costs better than many light ones.
\item The one-off \code{optimcon()} Constant broadcast scales with the
  size of drifts and operators; per-evaluation traffic does not. Sweep
  drivers that re-run \code{optimcon()} in a loop pay the broadcast each
  time --- restructure so that composition changes are outside the hot
  loop where possible.
\item Diagnostics cost real time every evaluation: switch
  \code{control.plotting} off for production runs and use
  \code{robustness} plots only while exploring.
\item On a single case, pool workers are idle by design; large
  single-system problems should invest in threaded BLAS or GPU operators
  instead.
\end{itemize}

%===============================================================
\section{Complete function and option reference}
\label{sec:reference}
%===============================================================

\subsection{\code{optimcon()} options: required fields}

\begin{center}
\begin{tabular}{lp{4.2cm}p{7.2cm}}
\toprule
field & type & meaning and constraints\\
\midrule
\code{drifts} & cell of cells of matrices & drift generators, outer
index over the ensemble, inner over time; inner length 1 or
\code{pulse\_ntpts}; all members equal inner length; dimensions must
match the control operators; Hermitian where unitarity is required\\
\code{operators} & cell of matrices & control operators
$\mathbf{C}_k$; cleaned up at absorption\\
\code{isotopes} & cell of strings & isotopes addressed by the control
channels; each must be present in the spin system\\
\code{channels} & vector of positive integers & index into
\code{isotopes} for each control operator, one element per operator\\
\code{rho\_init} & cell & initial states: column vectors
(\code{sphten-liouv}, \code{zeeman-liouv}, \code{zeeman-wavef}) or
square matrices (\code{zeeman-hilb})\\
\code{rho\_targ} & cell & target states, same count and kind as
\code{rho\_init}\\
\code{pwr\_levels} & row vector & nominal control powers, rad/s;
finite, positive; length $>1$ creates the power ensemble\\
\code{pulse\_dt} & row vector & interval durations, seconds; positive;
defines \code{pulse\_nsteps}, \code{pulse\_dur}, and
\code{pulse\_ntpts} (equal to \code{pulse\_nsteps}, or plus one for
the trapezium integrator)\\
\bottomrule
\end{tabular}
\end{center}

\subsection{\code{optimcon()} options: optional fields}

\begin{longtable}{lp{2.6cm}p{7.6cm}}
\toprule
field & default & meaning and constraints\\
\midrule
\endhead
\code{fidelity} & \code{'real'} & fidelity functional:
\code{'real'}, \code{'imag'}, or \code{'square'},
Eq.~\eqref{eq:fidelities}\\
\code{integrator} & \code{'rectangle'} & \code{'rectangle'}
(piecewise-constant) or \code{'trapezium'} (piecewise-linear)\\
\code{method} & \code{'lbfgs'} & optimiser: \code{'lbfgs'},
\code{'rbfgs'}, \code{'newton'}, \code{'goodwin'}; the last two refuse
the trapezium integrator and distortions; \code{'goodwin'} requires
Hermitian generators\\
\code{steady} & \code{false} & stroboscopic steady-state mode
(Section~\ref{sec:steady}); many restrictions apply\\
\code{prefix}, \code{suffix} & none & function handles
\code{rho=f(spin\_system,drift,rho)} applied before the pulse (initial
state) and to the target in reverse time\\
\code{dead\_time} & \code{0} & receiver dead time, seconds,
non-negative; target evolved backwards under the last drift\\
\code{offsets} + \code{off\_ops} & none & offset ensembles: cell of row
vectors (Hz) and matching cell of offset operators; both together;
same channel count; operator dimensions must match controls\\
\code{bsiegert} & \code{false} & logical scalar; enables
Bloch--Siegert corrections (Section~\ref{sec:bsiegert}): quadratic
response operators added to the slice generators and to the gradient;
requires first-order optimisers, the rectangle integrator, and
unit-quadrature control operators\\
\code{amplitudes} & --- & fixed amplitude profile for
\code{grape\_phase}; validated there\\
\code{freeze} & none & logical waveform-shaped mask of elements
excluded from optimisation; may not freeze everything; incompatible
with \code{basis}, which mixes time points\\
\code{max\_iter} & \code{100} & iteration cap, non-negative integer;
zero evaluates and plots once\\
\code{tol\_x} & \code{1e-3} & termination on step 1-norm\\
\code{tol\_g} & \code{1e-6} & termination on gradient 2-norm\\
\code{n\_grads} & \code{50} & LBFGS/RBFGS history depth, integer
$\geq2$\\
\code{phase\_cycle} & none & cycle lines
$[\varphi_{\mathrm{init}},\varphi_{1..K/2},\varphi_{\mathrm{targ}}]$,
radians; requires an even control count; $K/2+2$ columns\\
\code{basis} & none & orthonormal waveform basis, rows are basis
functions, \code{pulse\_ntpts} columns; optimisation runs in
coefficient space; incompatible with \code{freeze}\\
\code{keyholes} & none & cell of \code{pulse\_ntpts} entries, each
empty or a linear idempotent function handle; rectangle integrator
only\\
\code{penalties} + \code{p\_weights} & \code{\{'SNS'\}}, \code{100} &
penalty types from \code{'none'}, \code{'NS'}, \code{'SNS'},
\code{'SNSA'}, \code{'DNS'} with non-negative weights, both together,
equal counts\\
\code{u\_bound}, \code{l\_bound} & \code{+1}, \code{-1} & SNS
corridor in fractions of the power level; scalars or waveform-shaped
arrays; arrays are incompatible with SNSA and with amplitude-profile
optimisations\\
\code{ens\_corrs} & \code{\{\}} & ensemble correlations from
\code{'rho\_ens'}, \code{'rho\_drift'}, \code{'power\_drift'}\\
\code{budget} & \code{Inf} & ensemble budget: integer member count
($>1$) or fraction ($\leq1$); reproducible random subset\\
\code{plotting} & \code{\{\}} & diagnostic panels
(Section~\ref{sec:plotting}); expensive\\
\code{distplot} & \code{\{\}} & distortion chain applied to the
plotted waveform only\\
\code{distortion} & none & $E\times S$ cell of distortion function
handles with Jacobians; forces \code{lbfgs}\\
\code{traj\_opts} & \code{\{\}} & \code{\{'average'\}} returns the
ensemble-averaged trajectory\\
\code{checkpoint} & none & checkpoint file name in the scratch
directory; saves after every accepted step, restarts if present\\
\code{video\_file} & none & MJPEG video of the diagnostic figure;
requires plotting\\
\code{parameters} & none & free-form user data, absorbed without
parsing; visible to callbacks everywhere\\
\bottomrule
\end{longtable}

Unrecognised fields are reported by name and refused --- a typo cannot
silently become a no-op.

\subsection{Fields created by \code{optimcon()}}

For sweep drivers and callback authors, the frozen structure contains,
besides the absorbed options: \code{ncontrols} and \code{ndrifts}
(counts), \code{pulse\_dur}, \code{pulse\_nsteps}, \code{pulse\_ntpts}
(timing), \code{catalog} and
\code{ens\_sizes} (the ensemble case catalog and dimension sizes),
\code{invariants} (the worker-resident frozen problem,
a \code{parallel.pool.Constant}), \code{frozen\_fields} (the names of
the worker-resident invariant fields), the line search and regularisation
constants (\code{ls\_c1}, \code{ls\_c2}, \code{ls\_tau1},
\code{ls\_tau2}, \code{ls\_tau3}, \code{reg\_max\_iter},
\code{reg\_alpha}, \code{reg\_phi}, \code{reg\_max\_cond}), and the
parsed forms of every option above. The fields named in
\code{frozen\_fields} --- \code{drifts}, \code{operators},
\code{off\_ops}, and, when present, the trapezium commutator tables
\code{cc\_comm} and \code{cc\_comm\_idx} and the Bloch--Siegert
response operators \code{resp\_ops} --- are \emph{absent} from the
returned structure: they live on the workers.

\subsection{Public functions}

\begin{description}
\item[\code{spin\_system=optimcon(spin\_system,control)}] --- parses,
  validates, and freezes the optimal control problem; builds the
  ensemble catalog; publishes the worker-resident invariants.
\item[\code{[x,data]=fmaxnewton(spin\_system,cost\_function,guess)}]
  --- maximises the objective from \code{guess}; returns the optimal
  waveform and a diagnostics structure (\code{data.count.*} counters,
  \code{data.fx\_sep\_pen} fidelity-and-penalties vector of the last
  evaluation, \code{data.traj\_data} trajectory data).
\item[\code{[traj,fid,grad,hess]=grape\_xy(waveform,spin\_system)}] ---
  Cartesian cost function: ensemble fidelity with penalties separated
  (\code{fid(1)} fidelity, \code{fid(2:end)} weighted penalties;
  gradients and Hessians stacked accordingly in the third dimension).
\item[\code{[traj,fid,grad,hess]=grape\_phase(phi,spin\_system)}] ---
  phase-modulated cost function over a fixed amplitude profile.
\item[\code{[traj,fid,df\_du]=grape\_curv(u,u2x,dx\_du,spin\_system)}]
  --- curvilinear-coordinate cost function with user Jacobians.
\item[\code{[traj,fid,grad]=grape\_coop(phi,spin\_system)}] ---
  cooperative pulse pairs; the two phase profiles are stacked
  vertically.
\item[\code{[traj,fid,grad,hess]=ensemble(waveform,spin\_system)}] ---
  the parallel ensemble average of the bare fidelity (no penalties);
  called by the wrappers, callable directly.
\item[\code{tgrape()}] ---
  \code{[fid,grad]=tgrape(spin\_system,drift,controls,}
  \code{waveform,dt\_grid,time\_unit,rho\_init,rho\_targ)}:
  fidelity and its gradient with respect to interval durations at a
  fixed waveform.
\item[\code{[ops,coefs]=bloch\_siegert(spin\_system,ops,coefs)}] ---
  appends the Bloch--Siegert response operator of each control channel
  as a virtual channel whose coefficient vector is the square of that
  channel's amplitude vector; the outputs feed the piecewise-constant
  methods of \code{shaped\_pulse\_xy()}.
\item[\code{[catalog,ens\_sizes]=ens\_catalog(control)}] --- the
  ensemble case catalog: one row per case, columns indexing the state
  pair, drift, power, offset combination, phase cycle line, and
  distortion row.
\item[\code{[p,pg,ph]=penalty(wf,type,fb,cb)}] --- penalty value,
  gradient, and Hessian for a waveform between floor and ceiling
  bounds.
\item[\code{ctrl\_trajan(spin\_system,waveform,traj,fids)}] ---
  diagnostic plotting; called internally, configured through
  \code{control.plotting}.
\end{description}

\subsection{Low-level and internal functions}

\code{grape\_liouv()} and \code{grape\_hilb()} are the serial GRAPE
kernels (trajectories, gradients, Hessians for one ensemble case); they
accept the client-side structure with explicit drift and control
arguments and are the right entry point for custom ensemble schemes and
for regression tests. \code{aux\_mat()} builds the trapezium auxiliary
matrices, Eqs.~\eqref{eq:dgeffl}--\eqref{eq:dgeffr}; \code{dirdiff()}
(kernel utilities) implements the $N$-block directional derivatives;
\code{objeval()} assembles the total objective; \code{lbfgs()},
\code{bfgs()}, \code{hessreg()}, \code{bracketing()},
\code{sectioning()}, and \code{cubic\_interp()} are the optimiser
internals; \code{fapt2sfo()} and \code{inst\_freq()} support the
diagnostics; \code{bss\_ops()} builds the Bloch--Siegert response
operators from the channel isotope list and the carrier frequencies.

\subsection{Compatibility matrix}

\begin{center}
\begin{tabular}{lcc}
\toprule
feature & \code{rectangle} & \code{trapezium}\\
\midrule
LBFGS / RBFGS            & yes & yes\\
Newton / Goodwin Hessians & yes & no\\
keyholes                 & yes & no\\
waveform distortions     & LBFGS only & LBFGS only\\
Bloch--Siegert corrections & LBFGS/RBFGS only & no\\
stroboscopic steady state & yes (see below) & no\\
phase cycles, offsets, powers, drifts ensembles & yes & yes\\
prefix / suffix / dead time & yes & yes (not steady)\\
waveform basis            & yes & yes\\
\bottomrule
\end{tabular}
\end{center}

Steady-state mode additionally requires \code{sphten-liouv}, a
first-order optimiser, a traceless target, and no
prefix/suffix/dead-time; Hilbert space (\code{zeeman-hilb}),
wavefunction formalism, and the Goodwin method require all generators
Hermitian; phase cycles require an even number of controls; SNSA and
amplitude-profile optimisations require scalar bounds; array bounds
require SNS.

%===============================================================
\begin{thebibliography}{10}
%===============================================================

\bibitem{khaneja2005}
N.~Khaneja, T.~Reiss, C.~Kehlet, T.~Schulte-Herbr\"uggen, S.J.~Glaser,
\emph{Optimal control of coupled spin dynamics: design of NMR pulse
sequences by gradient ascent algorithms},
J.~Magn.~Reson.\ \textbf{172} (2005) 296--305.

\bibitem{defouquieres2011}
P.~de~Fouquieres, S.G.~Schirmer, S.J.~Glaser, I.~Kuprov,
\emph{Second order gradient ascent pulse engineering},
J.~Magn.~Reson.\ \textbf{212} (2011) 412--417.

\bibitem{goodwin2015}
D.L.~Goodwin, I.~Kuprov,
\emph{Auxiliary matrix formalism for interaction representation
transformations, optimal control, and spin relaxation theories},
J.~Chem.~Phys.\ \textbf{143} (2015) 084113.

\bibitem{goodwin2016}
D.L.~Goodwin, I.~Kuprov,
\emph{Modified Newton--Raphson GRAPE methods for optimal control of
spin systems},
J.~Chem.~Phys.\ \textbf{144} (2016) 204107.

\bibitem{najfeld1995}
I.~Najfeld, T.F.~Havel,
\emph{Derivatives of the matrix exponential and their computation},
Adv.~Appl.~Math.\ \textbf{16} (1995) 321--375.

\bibitem{vanloan1978}
C.F.~Van~Loan,
\emph{Computing integrals involving the matrix exponential},
IEEE Trans.\ Autom.\ Control \textbf{23} (1978) 395--404.

\bibitem{rasulov2023}
U.~Rasulov, A.~Acharya, M.~Carravetta, G.~Mathies, I.~Kuprov,
\emph{Simulation and design of shaped pulses beyond the
piecewise-constant approximation},
J.~Magn.~Reson.\ \textbf{353} (2023) 107478.

\bibitem{iserles1999}
A.~Iserles, S.P.~N{\o}rsett,
\emph{On the solution of linear differential equations in Lie groups},
Phil.~Trans.~R.~Soc.~Lond.~A \textbf{357} (1999) 983--1019.

\bibitem{liu1989}
D.C.~Liu, J.~Nocedal,
\emph{On the limited memory BFGS method for large scale optimization},
Math.~Program.\ \textbf{45} (1989) 503--528.

\bibitem{banerjee1985}
A.~Banerjee, N.~Adams, J.~Simons, R.~Shepard,
\emph{Search for stationary points on surfaces},
J.~Phys.~Chem.\ \textbf{89} (1985) 52--57.

\bibitem{fletcher2000}
R.~Fletcher,
\emph{Practical Methods of Optimization}, 2nd ed., Wiley, 2000.

\bibitem{hogben2011}
H.J.~Hogben, M.~Krzystyniak, G.T.P.~Charnock, P.J.~Hore, I.~Kuprov,
\emph{Spinach --- a software library for simulation of spin dynamics
in large spin systems},
J.~Magn.~Reson.\ \textbf{208} (2011) 179--194.

\bibitem{kuprov2023}
I.~Kuprov,
\emph{Spin: From Basic Symmetries to Quantum Optimal Control},
Springer, 2023.

\end{thebibliography}

\end{document}
